\documentclass[a4paper,12pt]{article}

\usepackage[T1]{fontenc}
\usepackage[italian]{babel}
\usepackage{graphicx}
\usepackage{geometry}
\usepackage{tabularx}
\usepackage[table]{xcolor}
\usepackage[most]{tcolorbox}
\usepackage{fancyhdr}
\usepackage[hidelinks]{hyperref}

\newcommand{\groupmail}{alittlebyte.swe@gmail.com}
\newcommand{\grouplogo}{../../../.github/site-src/assets/logo_alittlebyte.png}
\newcommand{\groupsymbol}{../../../.github/site-src/assets/solo_logo.png}

\geometry{
    left=2.5cm,
    right=2.5cm,
    top=2.5cm,
    bottom=2.5cm,
    headheight=331pt,
    includeheadfoot
}

\pagestyle{fancy}
\fancyhf{}

\renewcommand{\headrulewidth}{0pt}
\fancyhead{}

\renewcommand{\footrulewidth}{0.4pt}
\fancyfoot[L]{\includegraphics[height=0.6cm]{\groupsymbol}\hspace{0.45em}aLittleByte}
\fancyfoot[C]{\thepage}
\fancyfoot[R]{Verbale  2026-03-16}

\fancypagestyle{firstpage}{
    \fancyhf{}
    \renewcommand{\headrulewidth}{0pt}
    \renewcommand{\footrulewidth}{0.4pt}
    \fancyhead[C]{%
        \makebox[\textwidth][c]{%
            \begin{tabular}{c}
                \includegraphics[height=10cm]{\grouplogo}\\[1ex]
                \small\texttt{\groupmail}
            \end{tabular}%
        }
    }
    \fancyfoot[L]{\includegraphics[height=0.6cm]{\groupsymbol}\hspace{0.45em}aLittleByte}
    \fancyfoot[C]{\thepage}
    \fancyfoot[R]{Verbale  2026-03-16}
}

\begin{document}

\thispagestyle{firstpage}

\vspace*{0.5cm}

\begin{center}
    {\LARGE \textbf{Verbale Interno - 2026-03-16}}
\end{center}

\vspace{1.5cm}

\begin{center}
    \renewcommand{\arraystretch}{1.3}
    \begin{tcolorbox}[
        width=0.92\textwidth,
        colback=gray!6,
        colframe=gray!45,
        boxrule=0.5pt,
        arc=2mm,
        boxsep=0pt,
        left=8pt, right=8pt, top=8pt, bottom=8pt
    ]
        \begin{tabularx}{\linewidth}{>{\bfseries}p{3.5cm} X}
            Gruppo      & aLittleByte (gruppo 19) \\
            Redattore   & Angelica Gastaldello \\
            Verificatori& Leonardo Soligo\\
            Destinatari & Prof. Tullio Vardanega, Prof. Riccardo Cardin \\
            Data        &  2026-03-16\\
        \end{tabularx}
    \end{tcolorbox}
\end{center}

\newpage

\newgeometry{left=2.5cm,right=2.5cm,top=2.5cm,bottom=2.5cm,includefoot}
\tableofcontents
\newpage

\section{Informazioni Generali}
\section*{Specifiche Documento}
\hrule
\vspace{1cm}

\begin{tabular}{>{\bfseries}p{5cm} | p{10cm}}
Luogo Riunione & Discord \\[6pt]

Orario (Inizio - Fine) & 14.30 -- 15:45 \\[6pt]

Redattore &
  A. Gastaldello\\[6pt]

Amministratore & 
A. Gastaldello\\[6pt]

Verificatore &
L. Soligo\\[6pt]

Partecipanti &
  A. Oloni\\ &
    V. Y. Kaneda Fernandes\\&
  L. Soligo\\&
  A. Gastaldello\\&
  E. Bellet\\&
  D. Belluz\\&
  A. Maule\\[6pt]
  
\end{tabular}

\section{Ordine del Giorno}
    \begin{itemize}
    \item Domande per incontro col referente del Capitolato n.5 e n.1.
    \item Ragionamento sulle risposte fornite dal referente del Capitolato n.6.
    \item Organizzazione interna e workflow su GitHub.
    \item Aggiornamento stato documenti.
    \end{itemize}

\section{Attività svolte}
\subsection{Raccolta domande per il referente del Capitolato n.5 e n.1}
A seguito di un'analisi individuale del Capitolato n.5 e n.1, ciascun membro del gruppo ha raccolto osservazioni, dubbi e quesiti emersi durante la lettura. I contributi sono stati successivamente consolidati in un elenco condiviso, che costituirà la base delle richieste da sottoporre ai referenti dei capitolati nel prossimo incontro.

\subsection{Analisi delle risposte del referente del Capitolato n.6}
Il gruppo ha discusso le risposte fornite dal referente del Capitolato n.6 con l'obiettivo di valutarne l'adeguatezza rispetto alle esigenze e alle competenze del team. Il confronto e` stato inoltre orientato a definire come integrare in modo coerente gli esiti del colloquio nella documentazione, individuando i punti da aggiornare e i relativi approfondimenti.

\subsection{Organizzazione interna e workflow su GitHub}
Nel corso della riunione il gruppo ha definito i seguenti indirizzi operativi:
\begin{itemize}
    \item \textbf{Strategia di branching:} adozione di un branch dedicato per ciascun documento o attività, con l'obiettivo di evitare l'integrazione in main di contenuti non ancora validati.
    \item \textbf{Gestione task (Kanban):} adozione delle board GitHub Projects per rappresentare in modo trasparente lo stato di avanzamento delle attività e facilitare il coordinamento interno.
\end{itemize}
Le decisioni sopra riportate risultano pienamente coerenti con il way of working del gruppo, in quanto rafforzano tracciabilità, qualità documentale e responsabilità condivisa, riducendo al contempo il rischio di errori e conflitti nelle fasi di integrazione.

\subsection{Pianificazione dei prossimi passi su documentazione}
Il gruppo ha suddiviso i prossimi passi tra attività di scrittura e attività di revisione dei documenti. Per ogni documento sono stati assegnati almeno un redattore e un verificatore, così da garantire avanzamento ordinato e controllo della qualità.

\section{To-Do List}
\begin{table}[h]
    \centering
    \begin{tabular}{|p{4.2cm}|p{5cm}|l|}
    \hline
        \rowcolor{lightgray}\textit{\textbf{Persona Incaricata}}  & \textbf{Task Assegnata} &\textbf{Link Issue} \\
    \hline
        \textit{L. Soligo} & Creazione Kanban & \href{https://github.com/aLittleByte-19/Documentazione/issues/21}{\#21}\\
    \hline
        \textit{A. Gastaldello} & Redigere Verbale Interno 2026-03-16 & \href{https://github.com/aLittleByte-19/Documentazione/issues/22}{\#22}\\
    \hline
        \textit{A. Gastaldello e D. Belluz} & Aggiornare Analisi dei Capitolati dopo i colloqui & \href{https://github.com/aLittleByte-19/Documentazione/issues/23}{\#23}\\
    \hline
        \textit{A. Oloni, V. Y. Kaneda Fernandes, E. Bellet e A. Maule} & Approvazione documenti & \href{https://github.com/aLittleByte-19/Documentazione/issues/24}{\#24}, \href{https://github.com/aLittleByte-19/Documentazione/issues/25}{\#25}\\
    \hline
    \end{tabular}
    \label{tab:todo}
\end{table}

\end{document}